
\def\PUDcodigo{ECA223}

\if\COMPILAR0
   \def\PUDcodacad{04507.36}
\else
   \def\PUDcodacad{04506.36}
\fi

\def\PUDdisciplina{Máquinas Elétricas}

\def\PUDsemestre{S7}

\def\PUDhoras{80}

%NOVOS ---------------------------------------------------
\def\PUDhorasTeoria{80}
\def\PUDhorasPratica{0}

\if\COMPILAR0
    \def\PUDinterECA{\hyperref[PUD:ACE012]{Física 3 (04507.18)}
    
    \hyperref[PUD:ECA215]{Circuitos Elétricos II (04507.24)}
    
    \hyperref[PUD:ECA224]{Acionamento de Máquinas (04507.41)}
    }
\else
    \def\PUDinterECA{\hyperref[PUD:ACE012]{Física 3 (04506.18)}}
\fi

\def\PUDatuaisECA{---}

\def\PUDrecursos{Projetor multimídia; Quadro branco e pincel.}

%----------------------------------------------------------

\def\PUDcreditos{4}

\def\PUDprerequisito{ \hyperref[PUD:ACE012]{ACE012} }

\if\COMPILAR0
    \def\PUDpreacad{ \hyperref[PUD:ACE012]{Física 3 (04507.18)} }
\else
    \def\PUDpreacad{ \hyperref[PUD:ACE012]{Física 3 (04506.18)} }
\fi


\def\PUDobjetivos{Compreender os circuitos magnéticos.  Compreender o funcionamento dos transformadores monofásicos e trifásicos.  Compreender a conversão eletromecânica de energia.  Compreender o funcionamento das principais máquinas elétricas rotativas.}

\def\PUDementa{Circuitos magnéticos;  Transformadores;  Conversão eletromecânica de energia;  Máquinas de indução;  Máquinas síncronas;  %Máquinas de corrente contínua.
}

\def\PUDconteudo{ & Circuitos magnéticos.  Fluxo concatenado, Indução e Energia.  Propriedades dos materiais magnéticos.  Excitação CA. 
\\ 
 & Transformadores.  Circuito equivalente.  Transformadores de múltipos enrolamentos.  Transformadores especiais. 
\\ 
 & Forças e conjugados.  FMM de enrolamentos distribuídos.  Campos em máquinas rotativas.  Campos girantes em máquinas.  Conjugado.  Fluxos dispersivos. 
\\ 
 & Máquinas síncronas.  Circuito equivalente.  Ângulo de carga. 
\\ 
 & Máquinas de indução.  Circuito equivalente.  Ensaios básicos.  Efeitos da resistência do rotor. 
%\\ 
 %& Máquinas de corrente contínua.  Reação de armadura.  Comutação e interpolos.  Enrolamentos compensadores. 
 }

\def\PUDmetodo{Aulas expositórias que podem ser teóricas e/ou práticas, onde as práticas no laboratório serão marcadas no decorrer da disciplina.}

\def\PUDavalia{A avaliação será feita com aplicação de provas teóricas e/ou práticas, além da possibilidade de inclusão de trabalhos e seminários no decorrer da disciplina.}

\def\PUDbibbasica{ & \href{www.biblioteca.ifce.edu.br/index.asp?codigo_sophia=23684}{Fitzgerald}, A. E; Kingley Jr, Charles; Umans, Stephen D.  Máquinas Elétricas: com introdução à eletrônica de potência.  6º ed.  Porto Alegre, RS: Bookman, 2008.  648p.  ISBN: 9788560031047.
\\ 
 & \href{www.biblioteca.ifce.edu.br/index.asp?codigo_sophia=24384}{Kosow}, Irwing L.  Máquinas Elétricas e Transformadores.  15º ed.  São Paulo, SP: Globo, 2008.  667p.  ISBN: 8525002305.
\\ 
 & \href{www.biblioteca.ifce.edu.br/index.asp?codigo_sophia=24148}{Carvalho}, Geraldo.  Máquinas elétricas, teoria e ensaios.  4º ed.  São Paulo, SP: Érica, 2011.  263p.  ISBN: 9788536501260. }

\def\PUDbibcomplementar{ & \href{www.biblioteca.ifce.edu.br/index.asp?codigo_sophia=65036}{Filippo Filho}, Guilherme.  Motor de Indução: princípios de funcionamento, características operacionais, aplicações, acionamentos e comandos.  2º ed.  São Paulo, SP: Érica, 2016.  294p.  ISBN: 9788536504483.
\\ 
 & \href{www.biblioteca.ifce.edu.br/index.asp?codigo_sophia=24450}{Toro}, Vicent Del.  Fundamentos de máquinas elétricas.  Rio de Janeiro, RJ: LTC, 1999.  ISBN: 8521611846.
\\ 
 & \href{www.biblioteca.ifce.edu.br/index.asp?codigo_sophia=44617}{Rezek}, Angelo J.  Fundamentos básicos de máquinas elétricas: teoria e ensaios.  Rio de Janeiro, RJ: Synergia, 2011.  123p.  ISBN: 9788561325695.
\\
 & \href{www.biblioteca.ifce.edu.br/index.asp?codigo_sophia=40390}{Stephan}, Richard Magdalena.  Acionamento, comando e controle de máquinas elétricas.  Rio de Janeiro, RJ: Ciência Moderna, 2013.  230p.  ISBN: 9788539903542.
\\
 & \href{www.biblioteca.ifce.edu.br/index.asp?codigo_sophia=40033}{Simone}, Gilio Aluisio.  Conversão eletromecânica de energia: uma introdução ao estudo.  São Paulo, SP: Érica, 2014.  324p.  ISBN: 9788571946033. }

\def\PUDautor{Adriano Pereira}

\def\PUDdata{2013-04-28}

\def\PUDrevisao{3}

\def\PUDrevisor{Adriano Holanda Pereira}

\def\PUDdatarevisao{2019-05-15}

\PUDcorpo
